\chapter{Tableaux de calibration}
\label{ann:calibration}
% BUDGET : 2,5 pages. Version 2026-09-02 : l'annexe ne contient plus que des tableaux. Le
% balayage du plancher de prairie et l'argument exogene du cheptel, jusqu'ici en prose, sont
% devenus deux tables.
% Tous les nombres sont des macros lues dans les recap.json (build_chiffres.py), sauf le
% tableau du balayage plantain, saisi depuis CAL-76.

\section{\glsentryshort{pad} par culture, aux trois barreaux}

\begin{table}[H]
  \centering
  \caption{Surfaces observées et simulées par groupe \gls{rpg} (ha) et \gls{pad}
  correspondant (\%), aux trois barreaux de calibration. Le symbole $\bullet$ marque les
  cultures sous le seuil de \SI{\mesArtPad}{\percent} de \textcite{chopin2015}.
  \SI{\obsSurfaceIrreproductible}{\ha} du territoire ne sont reproductibles par aucune
  variante fine, ce qui impose un plancher d'écart de \num{\obsPlancherPad} point avant toute
  optimisation.}
  \label{tab:ann-pad-cultures}
  \scriptsize
  \begin{tabular}{@{}lrrrrrrr@{}}
    \toprule
    & & \multicolumn{2}{c}{\textbf{Parité \glsentryshort{gams}}}
      & \multicolumn{2}{c}{\textbf{+ plafond plantain}}
      & \multicolumn{2}{c}{\textbf{+ plancher prairie}} \\
    \cmidrule(lr){3-4}\cmidrule(lr){5-6}\cmidrule(lr){7-8}
    \textbf{Groupe} & \textbf{observé} & sim. & \gls{pad} & sim. & \gls{pad}
      & sim. & \gls{pad} \\
    \midrule
    Canne à sucre (CS) & \num{\obsHaCS} & \num{\simPariteCS} & \num{\padPariteCS}
      & \num{\simBcSeulCS} & \num{\padBcSeulCS} & \num{\simRetenuCS}
      & \num{\padRetenuCS}~$\bullet$ \\
    Prairie (PN) & \num{\obsHaPN} & \num{\simParitePN} & \num{\padParitePN}
      & \num{\simBcSeulPN} & \num{\padBcSeulPN} & \num{\simRetenuPN}
      & \num{\padRetenuPN}~$\bullet$ \\
    Banane export (BA) & \num{\obsHaBA} & \num{\simPariteBA} & \num{\padPariteBA}
      & \num{\simBcSeulBA} & \num{\padBcSeulBA} & \num{\simRetenuBA}
      & \num{\padRetenuBA}~$\bullet$ \\
    Maraîchage (MA) & \num{\obsHaMA} & \num{\simPariteMA} & \num{\padPariteMA}
      & \num{\simBcSeulMA} & \num{\padBcSeulMA} & \num{\simRetenuMA}
      & \num{\padRetenuMA}~$\bullet$ \\
    Jachère (JA) & \num{\obsHaJA} & \num{\simPariteJA} & \num{\padPariteJA}
      & \num{\simBcSeulJA} & \num{\padBcSeulJA} & \num{\simRetenuJA} & \num{\padRetenuJA} \\
    Plantain (BC) & \num{\obsHaBC} & \num{\simPariteBC} & \num{\padPariteBC}
      & \num{\simBcSeulBC} & \num{\padBcSeulBC} & \num{\simRetenuBC} & \num{\padRetenuBC} \\
    Igname (IG) & \num{\obsHaIG} & \num{\simPariteIG} & \num{\padPariteIG}
      & \num{\simBcSeulIG} & \num{\padBcSeulIG} & \num{\simRetenuIG} & \num{\padRetenuIG} \\
    Ananas (AN) & \num{\obsHaAN} & \num{\simPariteAN} & \num{\padPariteAN}
      & \num{\simBcSeulAN} & \num{\padBcSeulAN} & \num{\simRetenuAN} & \num{\padRetenuAN} \\
    Melon (ME) & \num{\obsHaME} & \num{\simPariteME} & \num{\padPariteME}
      & \num{\simBcSeulME} & \num{\padBcSeulME} & \num{\simRetenuME} & \num{\padRetenuME} \\
    Agrumes (AG) & \num{\obsHaAG} & \num{\simPariteAG} & \num{\padPariteAG}
      & \num{\simBcSeulAG} & \num{\padBcSeulAG} & \num{\simRetenuAG} & \num{\padRetenuAG} \\
    Vergers (VE) & \num{\obsHaVE} & \num{\simPariteVE} & \num{\padPariteVE}
      & \num{\simBcSeulVE} & \num{\padBcSeulVE} & \num{\simRetenuVE} & \num{\padRetenuVE} \\
    \midrule
    cultures sous le seuil & & & \num{\calibPariteCulturesOk}
      & & \num{\calibBcSeulCulturesOk} & & \num{\calibRetenuCulturesOk} \\
    \bottomrule
  \end{tabular}
\end{table}

\section{Les autres échelles, et la comparaison à l'article}

\begin{table}[H]
  \centering
  \caption{Les trois autres échelles d'évaluation, à la parité et sous la calibration retenue.
  Le seuil sous-régional et le seuil par ferme valent \SI{\mesArtPadZone}{\percent} et non
  \SI{\mesArtPad}{\percent}.}
  \label{tab:ann-echelles}
  \scriptsize
  \setlength{\tabcolsep}{4pt}
  \begin{tabularx}{\textwidth}{@{}p{28mm}rrX@{}}
    \toprule
    \textbf{Échelle} & \textbf{parité} & \textbf{retenu} & \textbf{Lecture} \\
    \midrule
    Sous-régions sous le seuil & & \num{\calibRetenuRegionsOk}/\num{\calibRetenuRegionsTotal} &
      les reçues tiennent entre \SI{\calibRetenuRegionsPadMin}{\percent} et
      \SI{\calibRetenuRegionsPadOkMax}{\percent}. Les deux échecs sont le Sud-Est et le
      Sud-Ouest de Basse-Terre, à \SI{\calibRetenuRegionsPadHorsMin}{\percent} et
      \SI{\calibRetenuRegionsPadMax}{\percent}. L'article échoue sur la même sous-région, le
      Sud-Est de Basse-Terre, sa seule à descendre à \SI{56}{\percent} de surface correctement
      simulée \\
    Fermes sous le seuil (\%) & &
      \num{\calibRetenuFermesOkPct} &
      \num{\calibRetenuFermesOk} fermes sur \num{\calibRetenuFermesTotal} \\
    \gls{pad} par ferme, médiane (\%) & \num{\calibPariteFermesPadMediane} &
      \num{\calibRetenuFermesPadMediane} &
      plus d'une ferme sur deux est reproduite exactement \\
    \gls{pad} par ferme, moyenne (\%) & \num{\calibPariteFermesPadMoyenne} &
      \num{\calibRetenuFermesPadMoyenne} &
      portée par une minorité de fermes très mal simulées. La lire comme un écart uniforme
      sous-estimerait largement ce que le modèle reproduit \\
    Types reproduits (\%) & \num{\calibPariteTypes} & \num{\calibRetenuTypes} &
      matrice de confusion des huit types d'exploitation \\
    Parcelles justes (\%) & & \num{\calibRetenuParcelles} &
      pour \SI{\calibRetenuSurface}{\percent} de la surface. Le modèle se trompe donc
      davantage sur les petites parcelles \\
    \midrule
    Marie-Galante & & \num{\calibRetenuRegionsPadMin} &
      passée du pire au meilleur rang. Sa prairie tombait de \SI{1357}{\ha} à \SI{228}{\ha}
      avant le plancher, soit \SI{83}{\percent} de perte sur une île entière \\
    \bottomrule
  \end{tabularx}
\end{table}

\begin{table}[H]
  \centering
  \caption{Les six corrections à appliquer avant de rapprocher ces chiffres de ceux de
  \textcite{chopin2015}.}
  \label{tab:ann-corrections}
  \scriptsize
  \setlength{\tabcolsep}{4pt}
  \begin{tabularx}{\textwidth}{@{}p{32mm}X@{}}
    \toprule
    \textbf{Point} & \textbf{Ce qu'il impose} \\
    \midrule
    Le run de parité n'est pas le modèle validé par l'article &
      l'article valide un modèle qui contient déjà le plafond de plantain, son équation~6.
      Notre parité le désactive pour mesurer ce qu'il apporte, et c'est la cause dominante de
      l'écart, le plantain y montant à \SI{\simPariteBC}{\ha} contre \SI{\obsHaBC}{\ha} \\
    L'article ne publie aucun \gls{pad} agrégé &
      la seule métrique comparable est le nombre de cultures sous le seuil, celle sur laquelle
      notre modèle est le plus faible. Ses dix usages fusionnent vergers et agrumes, les deux
      postes que nous ratons le plus \\
    Ses seuils sont deux et non un &
      \SI{\mesArtPad}{\percent} au territoire, \SI{\mesArtPadZone}{\percent} en sous-région et
      à la ferme. Comparer nos chiffres de ferme au premier serait gratuitement défavorable \\
    Les tables parcellaires n'ont pas le même dénominateur &
      la sienne compte les parcelles non cultivées, la nôtre non. Quatre points d'écart de
      définition \\
    Son année de base est 2010 &
      \num{\mesArtParcellesN} parcelles et \num{\mesArtFermes} exploitations contre
      \num{\obsParcelles} et \num{\obsExploitations} ici, soit \SI{13}{\percent}
      d'exploitations en moins à nombre de parcelles quasi identique. Les coefficients
      d'aversion de sa Table~2 étant calibrés sur ces fermes, les appliquer à 2017 est un
      transfert hors échantillon \\
    L'article ne publie aucune table par exploitation &
      notre quatrième échelle n'a pas de référence et ne se lit qu'en absolu \\
    \bottomrule
  \end{tabularx}
\end{table}

\section{Les balayages de seuil}

\begin{table}[H]
  \centering
  \caption{Balayage du plafond de plantain, sept seuils. Le \gls{pad} forme un palier plat de
  \SI{4650}{\tonne} à \SI{9240}{\tonne}, soit un facteur deux sur le seuil pour un demi-point
  d'écart. Le plantain simulé reste au-dessus des \SI{\obsHaBC}{\ha} observés à tous les
  seuils.}
  \label{tab:ann-balayage-plantain}
  \small
  \setlength{\tabcolsep}{4pt}
  \begin{tabular}{@{}lrrrrrrr@{}}
    \toprule
    \textbf{Seuil (t)} & \num{4650} & \num{6440} & \num{9240} & \num{15000} & \num{25000}
      & \num{40000} & désactivé \\
    \midrule
    \gls{pad} (\%) & \num{31.7} & \num{31.7} & \num{31.2} & \num{34.4} & \num{36.4}
      & \num{38.1} & \num{48.4} \\
    Types (\%) & \num{66.8} & \num{67.0} & \num{67.1} & \num{65.8} & \num{65.6} & \num{65.4}
      & \num{64.0} \\
    Surface (\%) & \num{69.0} & \num{69.0} & \num{69.3} & \num{68.6} & \num{68.2} & \num{67.7}
      & \num{64.3} \\
    Plantain (ha) & \num{179} & \num{247} & \num{355} & \num{577} & \num{960} & \num{1527}
      & \num{2875} \\
    \bottomrule
  \end{tabular}
\end{table}

\begin{table}[H]
  \centering
  \caption{Balayage du plancher de prairie, de la valeur retenue jusqu'à la statistique
  agricole. Les trois métriques que le plancher ne touche pas restent au-dessus des seuils
  publiés sur toute la plage, tandis que le \gls{pad} hors prairie empire quand le seuil
  monte.}
  \label{tab:ann-balayage-prairie}
  \scriptsize
  \setlength{\tabcolsep}{4pt}
  \begin{tabularx}{\textwidth}{@{}p{54mm}rX@{}}
    \toprule
    \textbf{Mesure} & \textbf{Valeur} & \textbf{Lecture} \\
    \midrule
    Plancher retenu & \SI{\mesCalalouBovins}{\ha} & valeur du paramètre \gls{gams}
      \texttt{QUOTA\_PN\_PIQ\_MIN} \\
    Plancher le plus haut balayé & \SI{\mesPrairieAgreste}{\ha} & prairie de la statistique
      agricole \\
    Types reproduits au plancher le plus haut & \SI{\mesPrairieHautTypes}{\percent} & \\
    Parcelles au plancher le plus haut & \SI{\mesPrairieHautParcelles}{\percent} & \\
    Surface au plancher le plus haut & \SI{\mesPrairieHautSurface}{\percent} & \\
    \gls{pad} hors prairie à \SI{7218}{\ha} & \SI{\mesPadHorsPrairieBas}{\percent} & \\
    \gls{pad} hors prairie à \SI{\mesPrairieAgreste}{\ha} & \SI{\mesPadHorsPrairieHaut}{\percent}
      & le plancher prend alors sa terre à la canne, la culture la mieux restituée. Un seuil
      calé pour flatter le résultat ferait exactement l'inverse \\
    \bottomrule
  \end{tabularx}
\end{table}

\begin{table}[H]
  \centering
  \caption{L'argument exogène qui fonde le plancher de prairie, et la décomposition du déficit
  de prairie observé sans plancher. La charge animale se déduit du recensement de 2017 et non
  du modèle.}
  \label{tab:ann-cheptel}
  \scriptsize
  \setlength{\tabcolsep}{4pt}
  \begin{tabularx}{\textwidth}{@{}p{54mm}rX@{}}
    \toprule
    \textbf{Grandeur} & \textbf{Valeur} & \textbf{Lecture} \\
    \midrule
    Bovins recensés en 2017 & \num{\mesBovins} & \\
    \gls{ugb} correspondantes & \num{\mesUgb} & \\
    Prairie de la statistique agricole & \SI{\mesPrairieAgreste}{\ha} &
      \num{\mesUgbHaAgreste} \gls{ugb} par hectare, conforme aux recensements \\
    Prairie du modèle sans plancher & \SI{\simBcSeulPN}{\ha} &
      \num{\mesUgbHaSansPlancher} \gls{ugb} par hectare, agronomiquement impossible \\
    Prairie de notre jeu parcellaire & \SI{\mesPrairieRpg}{\ha} & \\
    Statistique ramenée à notre couverture de \SI{\mesCouvertureSau}{\percent} &
      \SI{\mesPrairieExogene}{\ha} & \\
    \midrule
    Parcelles observées en prairie non reconduites sans plancher & \SI{3489}{\ha} &
      \SI{62}{\percent} d'économie légitime, \SI{36}{\percent} de couplage avec le plafond de
      main d'\oe uvre, \SI{2}{\percent} de résidu \\
    Part du déficit tenant à un écart de marge de \SI{0.94}{\percent} &
      \SI{30}{\percent} & \\
    Poids de la prairie dans la déviation totale avant son plancher & \SI{42}{\percent} & \\
    \bottomrule
  \end{tabularx}
\end{table}

\section{Ce que le solveur n'explique pas}

\begin{table}[H]
  \centering
  \caption{Les cinq suspects instruits puis écartés, avec la mesure qui a tranché. Les deux
  dernières lignes sont des verdicts tenus puis démentis.}
  \label{tab:ann-solveur}
  \scriptsize
  \setlength{\tabcolsep}{4pt}
  \begin{tabularx}{\textwidth}{@{}p{34mm}X@{}}
    \toprule
    \textbf{Suspect} & \textbf{Ce qui l'écarte} \\
    \midrule
    L'arbitraire du solveur & le même problème résolu avec trois graines de branchement rend
      \SI{\mesGraineA}{\percent}, \SI{\mesGraineB}{\percent} et \SI{\mesGraineC}{\percent} de
      \gls{pad} territorial. L'arbitraire vaut donc \num{\mesBruitPad} point, et c'est ce
      chiffre qui fixe à \SI{5}{\percent} la tolérance du garde-fou des runs de référence \\
    La main d'\oe uvre & le plan réellement observé demande \SI{\mesMoObserve}{\mega\hour} pour
      un plafond de \SI{\mesMoPlafond}{\mega\hour}. Le modèle n'a pas d'excuse de ce côté \\
    Recalibrer l'aversion au risque & descend le \gls{pad} à \SI{39.9}{\percent} avec des
      coefficients économiquement absurdes. Surajustement \\
    Caler des plafonds de marché sur l'observé & descend à \SI{33.6}{\percent} et c'est
      circulaire par construction \\
    Incorporer l'élevage & c'est la bonne réponse, et c'est pourquoi elle est écartée du stage
      plutôt que bâclée. Elle demande une variable d'état que le modèle mono-période ne porte
      pas \\
    \midrule
    \enquote{Le plancher de prairie rend le \gls{milp} intraitable} & verdict daté d'une
      version du modèle à \num{900000} variables. Il n'a pas survécu au masque
      d'éligibilité \\
    \enquote{Le plancher de prairie est circulaire} & mon propre verdict, tenu jusqu'au jour
      où la donnée exogène de charge animale a été trouvée \\
    \bottomrule
  \end{tabularx}
\end{table}

\begin{table}[H]
  \centering
  \caption{Les deux réserves à appliquer à la lecture des tableaux de cette annexe.}
  \label{tab:ann-reserves-calib}
  \scriptsize
  \setlength{\tabcolsep}{4pt}
  \begin{tabularx}{\textwidth}{@{}p{34mm}X@{}}
    \toprule
    \textbf{Réserve} & \textbf{Contenu} \\
    \midrule
    Le balayage du plantain est daté & ses valeurs viennent du balayage du 27 juillet, avant
      l'activation du plancher de prairie. Leur niveau n'est donc pas comparable au
      tableau~\ref{tab:ann-pad-cultures}, seule leur variation le long du balayage est en
      cause \\
    Chaque indicateur observé est une fourchette & l'assolement 2017 n'étant connu qu'au niveau
      des douze groupes agrégés, l'indicateur observé n'est pas une valeur. La subvention, le
      revenu, l'azote et le travail simulés tombent à l'intérieur de cette fourchette, et le
      point central en sort parfois par le haut. L'écart au central ne démontre donc rien \\
    \bottomrule
  \end{tabularx}
\end{table}
